% Coauthor briefing based on the manuscript and recorded experiment status.
% Build: latexmk -pdf -interaction=nonstopmode -halt-on-error report.tex
\documentclass[11pt]{article}
\usepackage[a4paper,margin=1.8cm]{geometry}
\usepackage[T1]{fontenc}
\usepackage{helvet}
\renewcommand{\familydefault}{\sfdefault}
\usepackage{booktabs,tabularx,array,xcolor,enumitem}
\usepackage{needspace}
\usepackage[hidelinks]{hyperref}
\setlist{nosep,leftmargin=*,topsep=3pt}
\setlength{\parindent}{0pt}
\setlength{\parskip}{5pt}
\definecolor{accent}{HTML}{1F7F83}
\newcommand{\hd}[1]{\needspace{5\baselineskip}\vspace{7pt}{\large\bfseries\color{accent}#1}\par\vspace{2pt}}
\newcommand{\sub}[1]{\needspace{3\baselineskip}\vspace{3pt}{\bfseries #1}\par}
\newcolumntype{L}{>{\raggedright\arraybackslash}X}
\renewcommand{\arraystretch}{1.16}

\begin{document}

{\LARGE\bfseries VCF-RDFizer: coauthor discussion brief}\\[4pt]
{\small 1 October 2026 \textbullet\ Current manuscript and next decisions}

\hd{1. What is the paper's argument?}
\textbf{VCF-RDFizer turns genomic variant files into a shared graph with explicit links and access rules. These benefits must justify the cost of conversion.}

VCF is a file format for genomic variants. RDF represents those data as a graph;
each statement is a \emph{triple}. VCF-RDFizer adds shared variant identifiers,
links to external resources, and machine-readable access policies. SPARQL is the
language used to query that graph. Conventional VCF tools can answer the same
biological questions; the proposed benefit is making the links and rules explicit
and independently checkable.

\textbf{Meeting goal:} agree on this central claim, the experiments needed before
submission, and the detail that can move out of the main text.

\hd{2. What has been demonstrated?}
\sub{A question that needs both external knowledge and access rules}
\emph{Which participants carry a ClinVar-classified variant in the 81 genes on
the ACMG secondary-findings list, and which results may each requester receive?}

ClinVar supplies variant classifications. We use real public genomes but
\textbf{simulated consents}: some participants permit research, some permit only
clinical care, and some have withdrawn. A further rule restricts results in 28
cancer-predisposition genes to clinical care. A clinical laboratory, cardiovascular
study and biobank therefore receive different results.

Two independent implementations answer the same question: VCF-RDFizer with graph
queries and policies, and a conventional \texttt{bcftools} pipeline with Python
access-rule checks. \textbf{Their answers match for every requester in all three
setups.}

\begin{center}
\small
\begin{tabularx}{\textwidth}{@{}L r r r@{}}
\toprule
\textbf{Data tested} & \textbf{Clinical lab} & \textbf{Cardio study} & \textbf{Biobank}\\
\midrule
5 genomes, restricted to the 81 genes & 2{,}878 & 2{,}988 & 1{,}987\\
104 participants, the same gene regions & 63{,}529 & 92{,}870 & 48{,}207\\
1 whole genome (HG005) & 1{,}496 & 0 & 0\\
\bottomrule
\end{tabularx}
\end{center}
{\small Counts are participant--variant--gene matches, \textbf{not numbers of people}.
HG005 permits clinical care only; its whole-genome answer matches its gene-region
answer. The whole-genome graph contains 668 million triples before links.}

\sub{What this result means}
\begin{itemize}
\item Shared identifiers join the genomes to ClinVar even when files name the
same chromosome differently. The cohort also joins a separate frequency file to
find rare variants, again matching the conventional pipeline.
\item Each released dataset is checked against the access rules and against an
independent result derived from the source VCF. Agreement supports correct
implementation of these rules; it does not establish real-world consent validity.
\item The question includes \emph{all} ClinVar classifications, not only harmful
variants. The pathogenic/likely pathogenic subset is empty in the five-genome
setup and contains 8 carrier matches in the cohort.
\end{itemize}

\textbf{Effort is not yet a strong selling point.} The authored rules take 72 lines
in the RDF route versus 131 in the conventional route, but three of four tested
changes are similarly small in both. A rule on one named variant needs policy
data only in RDF and a small code change in the baseline. This is evidence of
flexibility, not a general productivity study.

\newpage
{\Large\bfseries Evidence and limits}

\hd{3. How well is the conversion checked?}
The base benchmark comprises 143 recorded cases across ten public files and controlled
record- and sample-count experiments. Validation compares RDF answers with the
original VCF and uses SHACL rules to check graph structure and values.

\begin{tabularx}{\textwidth}{@{}>{\raggedright\arraybackslash}p{0.25\textwidth} L@{}}
\toprule
\textbf{Check} & \textbf{Evidence and scope}\\
\midrule
13 paired query checks & 984 comparisons match; 4 are not applicable; none disagree. The full validation campaign covers fixtures and slices up to 17.1M triples.\\
\addlinespace
Injected faults & Queries detect 96/113; adding all SHACL profiles detects 113/113. The default profile's detection rate is not yet reported.\\
\addlinespace
Compressed files & All 152 HDT/COTTAS files decode with matching triple counts. Counts alone do not establish that every value is correct.\\
\bottomrule
\end{tabularx}

\textbf{A wording gap to fix.} The scale-retrieval runs also compare query answers
with the VCF at 657M triples, but record timing results rather than a complete
validation verdict. The paper should distinguish this query agreement from the
full validation pipeline; ``the whole genome received only count checks'' is too
broad. Agreement with a VCF also does not establish that its biological calls are
correct.

\hd{4. When is the extra cost worthwhile?}
\begin{tabularx}{\textwidth}{@{}>{\raggedright\arraybackslash}p{0.29\textwidth} L@{}}
\toprule
\textbf{Workload} & \textbf{Finding and implication}\\
\midrule
One question or one batch of summaries & Reading the VCF directly is faster once graph conversion and query-index setup are included.\\
\addlinespace
Repeated questions on an existing graph & QLever, the fastest tested graph engine, answers individual questions faster after indexing. Conversion and setup break even after about 44 questions on the 100,000-record test; this is workload-specific.\\
\addlinespace
Small genomic-region lookups & Coordinate-indexed VCF access (tabix) wins at 1\,kb windows; QLever wins at 10\,Mb. The fastest route depends on the question.\\
\addlinespace
Whole-genome conversion & One benchmark run took 16.1\,h, including 14.5\,h building HDT and COTTAS. Its total retained outputs occupied about 78 times the compressed VCF download.\\
\bottomrule
\end{tabularx}

\sub{What the representation choices mean}
\textbf{HDT and COTTAS are compressed RDF formats.} QLever builds its own index
from either, so the format barely affects its query time. Their main measured
role here is storage and transfer. Engines querying them directly did not
complete the tested 171M-triple workload.

\textbf{Sample profiles trade direct access for size.} Expanded representation
makes each sample value a graph resource. Condensed packs values into vectors:
smaller, but individual access needs decoding. At 2,504 samples, expanded has
432 times as many triples but 69 times the HDT bytes. Compare bytes as well as
triple counts.

\textbf{Do not mix whole-genome timings.} The 16.1-hour benchmark includes
compressed-format construction. The clinical use case's 53-minute conversion
uses another configuration and later software, with linking, release and indexing
costs reported separately. This is not a before/after speed-up.

\newpage
{\Large\bfseries What should we address before submission?}

\hd{5. Current state and immediate work}
The use case, performance comparisons and manuscript restructure are complete
in the current draft. The narrative is about 23 pages (32 with declarations and
references), plus a 9-page supplement. The earlier implementation notes still
list the following two experiments as pending; their results are not in the draft.

\begin{tabularx}{\textwidth}{@{}>{\raggedright\arraybackslash}p{0.29\textwidth} L@{}}
\toprule
\textbf{Priority} & \textbf{Action and reason}\\
\midrule
Test the default validation settings & Repeat fault injection with the default SHACL profile. The reported 113/113 detection uses all profiles, which are more expensive. Readers need the sensitivity of the settings they normally run.\\
\addlinespace
Check a heterogeneous real input & Run full paired validation on a 250,000-record genome slice from a different sequencing provider. This tests value preservation beyond the existing fixtures and reference-genome slice.\\
\addlinespace
Align the claims and evidence & Clarify validation coverage, label simulated consents and carrier-match counts, and keep setup costs beside speed claims. Describe controlled release as access-rule enforcement, not anonymization.\\
\addlinespace
Finish presentation and provenance & Aim for roughly 20 narrative pages. Complete funding, archive identifiers and release metadata. The base campaign used v3.1.0; the use case used v3.2.0 conversion and v3.3.0 plug-ins. Confirm the final release and automated-test status before calling it complete.\\
\bottomrule
\end{tabularx}

\hd{6. Decisions for the coauthors}
\begin{enumerate}[itemsep=8pt]
\item \textbf{Is the demonstrated benefit convincing?} Are shared identifiers,
inspectable policies and checked release sufficient with simulated consents, or
do we need an external user or data-access partner? Avoid claiming that fewer
lines establish usability.
\item \textbf{How much validation is enough for this paper?} Are the two targeted
runs above sufficient with a precise scope statement, or is complete whole-genome
validation required? The latter needs more work on the costly shape checks.
\item \textbf{Do we need a measured converter comparison?} The current comparison
with other converters is based on documentation. Decide whether to add a bounded
comparison on one shared input, rather than imply measured superiority.
\item \textbf{What belongs in the main story?} Keep conversion, validation and the
linked-data use case central; treat HDT/COTTAS as measured storage choices. Agree
whether the policy component belongs here or warrants a separate paper, and how
this submission relates to the companion vocabulary paper.
\item \textbf{What can wait?} Keep federation with other graph services,
alternative variant identifiers, whole-cohort conversion and sample-level access
control in multi-sample files as future work unless one is essential to the
paper's central claim. Choose which detailed tables move to the supplement.
\end{enumerate}

\vfill
{\footnotesize\color{gray}\raggedright
Reading map: manuscript Sections 4.1 (fidelity), 4.2 (use case), 4.3--4.4
(costs), and 5 (interpretation); Supplementary Sections S1--S5 (execution,
coverage and testing issues). Experiment status follows
\nolinkurl{tool-docs/workstream-a-implementation.md} and the archived scale-retrieval
reports.\par}

\end{document}
